\begin{table}[!ht]
\centering
\caption{Civilian Department Heterogeneity}
\label{tab:dept}
\footnotesize
\begin{adjustbox}{max width=\textwidth}
\begin{tabular}{lrrrr}
\toprule
Department & Pre (avg/mo) & Post (avg/mo) & \% change & DiD est. \\
\midrule
Commerce & 398 & 84 & -78.8 & -2.475 (0.115) \\
Agriculture & 1232 & 192 & -84.4 & -1.963 (0.115) \\
Other Agencies and Independent Organizations & 1546 & 286 & -81.5 & -1.340 (0.115) \\
Health and Human Services & 1035 & 280 & -72.9 & -0.957 (0.115) \\
Interior & 1326 & 368 & -72.2 & -0.913 (0.115) \\
National Aeronautics and Space Administration & 155 & 61 & -60.5 & -0.414 (0.115) \\
Transportation & 573 & 267 & -53.4 & -0.003 (0.115) \\
Veterans Affairs & 8452 & 4186 & -50.5 & 0.101 (0.115) \\
Justice & 992 & 671 & -32.3 & 0.375 (0.115) \\
Homeland Security & 1605 & 1259 & -21.6 & 0.518 (0.115) \\
Legislative Branch & 166 & 169 & 2.3 & 0.802 (0.115) \\
Judicial Branch & 53 & 59 & 11.6 & 0.865 (0.115) \\
\bottomrule
\end{tabular}
\end{adjustbox}
\begin{minipage}{0.92\textwidth}
\footnotesize\textit{Notes:} For each civilian department, we estimate a separate DiD with that department as the only treated unit and the four military departments as controls (two-way FE, log(1+vacancies)). Pre/Post columns report the average monthly vacancy count in Jan 2021--Jan 2025 vs Feb--Mar 2025. SEs in parentheses, clustered at department.
\end{minipage}
\end{table}
